% Section 4 — Key findings and results (part 1: 4.1-4.10)
\section{Key Findings and Results}
\label{sec:findings}

\subsection{Overall forecasting results}
\label{sec:overall-fc}
\Cref{fig:mase-board} tells the whole forecasting story in one panel. On M5, a
tight lead cluster --- LSTM (1.316), Moving Average (1.346), SES (1.351), TSB
(1.369) --- separates from ARIMA (1.505), the Naive family ($\approx$1.53),
and the smoothing failures DES (1.987) and TES (2.599). On Store demand the
spread is wider and more decisive: LSTM (0.978) stands alone below 1.0,
followed by SARIMA-full500 (1.055), Seasonal Naive (1.183), and Moving Average
(1.191), with Naive (1.971) and TES (2.111) far behind. Two remarks before the
detail: MASE below 1.0 means beating a seasonal-naive-on-training benchmark,
which only Store LSTM achieves; and the M5 lead cluster is so tight that
statistical machinery is needed to rank inside it honestly
(\cref{sec:spotlight}).

\begin{figure}[H]
\centering
\includegraphics[width=\textwidth]{F-NEW-01-combined-mase-leaderboard.png}
\caption{Forecast leaderboard (MASE, $m=7$, lower is better). M5: tight LSTM /
MA / SES / TSB cluster. Store: LSTM alone sub-1.0, SARIMA-full500 second.
Deterministic rebuild from final files (F-NEW-01).}
\label{fig:mase-board}
\end{figure}

\subsection{M5 forecasting results}
\label{sec:m5-fc}
\Cref{tab:m5fc} is the M5 board. LSTM leads on both MAE (0.989) and MASE
(1.316); Moving Average (0.9997 / 1.346) and SES (0.998 / 1.351) are within a
whisker, and TSB (1.010 / 1.369) completes the lead pack. SBA and Croston sit
just behind; ARIMA (1.132 / 1.505) trails the pack without disgrace; the Naive
pair ($\approx$1.19) anchors the middle; DES (1.514) and TES (1.922) fail
outright. Note the RMSE column's warning: SES's RMSE (2.00) exceeds
Croston's (2.11)? No --- read carefully: SES 2.00 vs.\ Croston 2.11 vs.\ LSTM
1.95 vs.\ Moving Average 1.30. Moving Average owns the best RMSE by far,
because averaging suppresses the large errors RMSE punishes quadratically ---
yet it loses MAE/MASE to LSTM. Different losses already crown different
champions before inventory enters the picture. sMAPE ($\approx$117--161\%)
confirms that percentage errors on near-zero demand are huge for everyone;
WAPE (per-series mean) ranks LSTM (0.77) first, matching the MAE story.

\begin{table}[H]
\centering
\caption{M5 forecast leaderboard ($n=112{,}000$ per model). Bold: column best.}
\label{tab:m5fc}
\begin{tabular}{lcccccc}
\toprule
Model & MAE & RMSE & MASE & RMSSE & WAPE & sMAPE (\%) \\
\midrule
\textbf{LSTM} & \textbf{0.99} & 1.95 & \textbf{1.32} & \textbf{0.85} & \textbf{0.77} & 144.2 \\
Moving Average & 1.00 & \textbf{1.30} & 1.35 & 0.86 & 1.31 & 143.5 \\
SES & 1.00 & 2.00 & 1.35 & 0.86 & 0.78 & 145.1 \\
TSB & 1.01 & 2.01 & 1.37 & 0.88 & 0.79 & 146.3 \\
SBA & 1.07 & 2.12 & 1.41 & 0.89 & 0.83 & 148.1 \\
Croston & 1.08 & 2.11 & 1.42 & 0.87 & 0.85 & 143.9 \\
ARIMA & 1.13 & 2.32 & 1.51 & 1.05 & 0.88 & 139.0 \\
Naive & 1.18 & 1.60 & 1.55 & 1.12 & 1.43 & 160.9 \\
Seasonal Naive & 1.19 & 1.67 & 1.53 & 1.14 & 1.45 & 152.7 \\
DES & 1.51 & 3.43 & 1.99 & 1.33 & 1.18 & 117.8 \\
TES & 1.92 & 4.39 & 2.60 & 1.85 & 1.50 & 116.0 \\
\bottomrule
\end{tabular}
\end{table}

\begin{figure}[H]
\centering
\includegraphics[width=0.95\textwidth]{01_comparison_MAE.png}
\caption{Full-ladder MAE comparison including LSTM (hardened run). The M5
cluster at $\approx$1.0 and Store's LSTM-led spread are both visible here;
exact boards are \cref{tab:m5fc,tab:storefc}. Source: notebook 08.}
\label{fig:mae-ladder}
\end{figure}
\begin{figure}[H]
\centering
\includegraphics[width=0.95\textwidth]{01_comparison_WAPE.png}
\caption{Full-ladder WAPE comparison. Note Store WAPE values are an order of
magnitude smaller (dense demand inflates the denominator) --- cross-dataset
WAPE levels are not comparable; ranks are. Source: notebook 08.}
\label{fig:wape-ladder}
\end{figure}

\subsection{Store forecasting results}
\label{sec:store-fc}
\Cref{tab:storefc} is the Store board, and it is less ambiguous. LSTM wins
every column that matters (MAE 7.87, MASE 0.978, RMSSE 0.971, WAPE 0.138);
SARIMA-full500 is a clear second (8.45 / 1.055 / 1.040 / 0.148);
Seasonal Naive (9.45) and Moving Average (9.71) form the next tier; SES (9.92)
follows; ARIMA (11.90), DES (13.63), Naive (16.29), TES (16.87) trail. Two
points deserve emphasis. First, seasonal structure pays on Store demand:
SARIMA's seasonal moving-average term cuts MAE from ARIMA's 11.90 to 8.45, and
even the humble Seasonal Naive beats SES. Second, the full500 SARIMA (MAE
8.45, MASE 1.055) differs modestly from the exploratory 100-series subset
(8.31 / $\approx$1.069): a population shift, documented, not hidden --- and a
reminder that subset results must never headline.

\begin{table}[H]
\centering
\caption{Store forecast leaderboard ($n=112{,}000$ per model; SARIMA full500). Bold: column best.}
\label{tab:storefc}
\begin{tabular}{lcccccc}
\toprule
Model & MAE & RMSE & MASE & RMSSE & WAPE & sMAPE (\%) \\
\midrule
\textbf{LSTM} & \textbf{7.87} & \textbf{10.42} & \textbf{0.98} & \textbf{0.97} & \textbf{0.14} & \textbf{15.3} \\
SARIMA & 8.45 & 11.09 & 1.05 & 1.04 & 0.15 & 16.9 \\
Seasonal Naive & 9.45 & 11.64 & 1.18 & 1.17 & 0.18 & 19.2 \\
Moving Average & 9.71 & 12.08 & 1.19 & 1.18 & 0.18 & 18.3 \\
SES & 9.92 & 13.41 & 1.21 & 1.21 & 0.17 & 18.6 \\
ARIMA & 11.90 & 15.98 & 1.45 & 1.43 & 0.21 & 23.0 \\
DES & 13.63 & 18.44 & 1.68 & 1.67 & 0.24 & 25.4 \\
Naive & 16.29 & 18.95 & 1.97 & 1.87 & 0.29 & 33.2 \\
TES & 16.87 & 23.39 & 2.11 & 2.16 & 0.29 & 33.7 \\
\bottomrule
\end{tabular}
\end{table}

\begin{figure}[H]
\centering
\includegraphics[width=0.9\textwidth]{F-NEW-02-store-mae-full500-sarima.png}
\caption{Store MAE with full-population SARIMA (8.45, $n=112{,}000$). This
replaces the stale subset-based panels. Deterministic rebuild (F-NEW-02).}
\label{fig:store-mae}
\end{figure}

\subsection{Statistical significance vs.\ practical significance}
\label{sec:sig-vs-pract}
With 112,000 paired errors per comparison, a microscope finds everything
significant. \Cref{tab:spotlight} therefore reports effect sizes alongside
$p$-values, and the lesson is taught by two rows. M5 LSTM vs.\ TSB:
$d_z \approx -0.031$, Holm-adjusted $p \approx 2\times10^{-10}$ ---
astronomically significant, practically tiny: the models differ by two
hundredths of an MAE unit. Store LSTM vs.\ Moving Average: $d_z \approx -0.29$,
$p \approx 0$ --- significant \emph{and} material, nearly two MAE units apart.
Same ``significant'' label, tenfold difference in meaning. A $p$-value answers
``are they different?''; only $d_z$, the rank-biserial $r$, and the bootstrap
interval answer ``should anyone care?'' \Cref{app:stats} formalizes the
machinery; the full 91-pair table is \cref{app:pairs}.

\begin{table}[H]
\centering
\caption{Spotlight pairwise comparisons (paired-error $d$ = A minus B; negative favors A).}
\label{tab:spotlight}
\begin{tabular}{lllccc}
\toprule
Dataset & A vs.\ B & $d_z$ & Holm $p$ & $r$ & 95\% CI \\
\midrule
M5 & LSTM vs.\ TSB & $-0.031$ & $2.0{\times}10^{-10}$ & $-0.020$ & $[-0.024,-0.017]$ \\
M5 & LSTM vs.\ MA & $-0.015$ & $1.8{\times}10^{-3}$ & $-0.010$ & $[-0.014,-0.007]$ \\
M5 & LSTM vs.\ SES & $-0.014$ & $5.1{\times}10^{-4}$ & $-0.011$ & $[-0.013,-0.005]$ \\
M5 & MA vs.\ SES & $+0.006$ & $1.7{\times}10^{-12}$\rlap{\textsuperscript{*}} & $-0.022$ & $[0.000,0.003]$ \\
Store & LSTM vs.\ MA & $-0.295$ & $\approx$0 & $-0.282$ & $[-1.868,-1.800]$ \\
Store & LSTM vs.\ SARIMA & $-0.118$ & $\approx$0 & $-0.128$ & $[-0.599,-0.548]$ \\
Store & LSTM vs.\ SES & $-0.336$ & $\approx$0 & $-0.311$ & $[-2.085,-2.014]$ \\
\bottomrule
\end{tabular}
\end{table}
\noindent\textsuperscript{*}Wilcoxon significant but Diebold--Mariano $p=0.056$
(raw) / $0.113$ (Holm): the two tests disagree on MA vs.\ SES --- correctly, the
pair is declared indistinguishable in practice.

\subsection{Pairwise statistical comparisons}
\label{sec:spotlight}
Across all 91 pairs, three patterns stand out. First, M5's lead cluster is
genuinely ordered (LSTM $>$ MA $\approx$ SES $>$ TSB, all pairwise significant)
but the gaps are operationally negligible --- the entire cluster spans 0.05
MASE. Second, Store gaps are an order of magnitude larger and every LSTM
pairing is both significant and material. Third, the MA--SES pair on M5 is the
honest null of the study: Wilcoxon says significant, Diebold--Mariano says not
($p=0.113$ Holm), the interval straddles zero --- two excellent simple methods
with nothing between them. That disagreement is a feature: it shows the
machinery can say ``too close to call'' even at $n=112{,}000$.

\subsection{Inventory evaluation framework}
\label{sec:inv-frame}
Recall the shared engine from \cref{sec:policy}: every forecast drives the
same daily-review order-up-to policy ($L=7$, 95\% target, $H=1$, $P=5$), lost
sales, cost $=$ holding $+$ stockout. Rankings below are total cost under that
policy; service level, average inventory, and reorder frequency diagnose
\emph{how} each cost was earned. Sensitivity to the policy itself is not a
footnote --- it is \cref{sec:robust,sec:sens}.

\subsection{M5 inventory results}
\label{sec:m5-inv}
\Cref{tab:m5inv} ranks M5 decision quality. LSTM wins (152.83), daylight ahead
of Moving Average (158.83) and SES (159.31); SBA (160.66), TSB (165.13), and
Croston (174.81) follow; Seasonal Naive (192.55), ARIMA (194.60), and Naive
(200.75) form the expensive middle; DES (268.80) and TES (346.67) are in a
different, worse league. The cost anatomy matters: SBA's cheap holding (128.39)
is eaten by stockouts (32.27); SES pairs the lowest stockout bill (21.00) with
the second-best service level (93.6\%); Croston buys the best service (93.9\%)
at middling cost. LSTM's margin ($\approx$6 over MA) is ten times its MAE
margin --- the policy amplifies small accuracy edges into real money when
demand is sparse.

\begin{table}[H]
\centering
\caption{M5 inventory ranking under the shared policy (cost units; lower better). Bold: best.}
\label{tab:m5inv}
\begin{tabular}{lccccc}
\toprule
Model & Total & Holding & Stockout & Service (\%) & Avg.\ inv.\ \\
\midrule
\textbf{LSTM} & \textbf{152.83} & 130.04 & 22.80 & 92.8 & 4.64 \\
Moving Average & 158.83 & 136.30 & 22.52 & 93.0 & 4.87 \\
SES & 159.31 & 138.32 & 21.00 & 93.6 & 4.94 \\
SBA & 160.66 & 128.39 & 32.27 & 90.9 & 4.59 \\
TSB & 165.13 & 143.54 & 21.59 & 93.2 & 5.13 \\
Croston & 174.81 & 154.54 & 20.28 & 93.9 & 5.52 \\
Seasonal Naive & 192.55 & 170.73 & 21.82 & 93.8 & 6.10 \\
ARIMA & 194.60 & 162.62 & 31.98 & 89.6 & 5.81 \\
Naive & 200.75 & 158.25 & 42.50 & 87.2 & 5.65 \\
DES & 268.80 & 235.35 & 33.46 & 89.9 & 8.41 \\
TES & 346.67 & 304.28 & 42.40 & 88.0 & 10.87 \\
\bottomrule
\end{tabular}
\end{table}

\subsection{Store inventory results}
\label{sec:store-inv}
\Cref{tab:storeinv} ranks Store decision quality --- and upends the forecast
board. Moving Average wins (2,084.50); SES (2,113.83) and Seasonal Naive
(2,128.15) complete a podium with no LSTM on it; LSTM (2,247.45) sits fourth,
SARIMA (2,282.31) fifth, ARIMA (2,320.87) sixth; Naive (2,700.37), DES
(2,779.64), TES (3,419.88) trail. The anatomy flips versus M5: Naive holds the
least stock (28.38) but pays a catastrophic stockout bill (1,905.74) at 43.2\%
service; DES/TES drown in inventory (76.07 / 88.77) yet still pay mid-range
stockouts --- overstocking without protection. Moving Average's virtue is
balance: moderate stock (45.58), the lowest stockout frequency (10.72), and no
extreme in either column. \textbf{This ranking is policy-contingent} ---
\cref{sec:robust} shows the crown moves --- and must never be quoted without
that caveat.

\begin{table}[H]
\centering
\caption{Store inventory ranking under the shared policy. Bold: best. Caution: fragile across policies (\cref{sec:robust}).}
\label{tab:storeinv}
\begin{tabular}{lccccc}
\toprule
Model & Total & Holding & Stockout & Service (\%) & Avg.\ inv.\ \\
\midrule
\textbf{Moving Average} & \textbf{2,084.50} & 1,276.10 & 808.40 & 61.7 & 45.58 \\
SES & 2,113.83 & 1,155.60 & 958.24 & 56.1 & 41.27 \\
Seasonal Naive & 2,128.15 & 1,281.05 & 847.10 & 62.0 & 45.75 \\
LSTM & 2,247.45 & 1,157.70 & 1,089.76 & 53.1 & 41.35 \\
SARIMA & 2,282.31 & 1,307.78 & 974.53 & 59.2 & 46.71 \\
ARIMA & 2,320.87 & 1,076.48 & 1,244.40 & 52.4 & 38.45 \\
Naive & 2,700.37 & 794.63 & 1,905.74 & 43.2 & 28.38 \\
DES & 2,779.64 & 2,130.09 & 649.55 & 73.5 & 76.07 \\
TES & 3,419.88 & 2,485.65 & 934.23 & 69.1 & 88.77 \\
\bottomrule
\end{tabular}
\end{table}

\begin{figure}[H]
\centering
\includegraphics[width=0.95\textwidth]{01_total_cost_comparison.png}
\caption{Total inventory cost by model and dataset under the shared policy.
M5: LSTM leads. Store: Moving Average leads, LSTM fourth. Source: unified
simulation (notebook inventory layer).}
\label{fig:totalcost}
\end{figure}
\begin{figure}[H]
\centering
\includegraphics[width=0.95\textwidth]{02_service_level_comparison.png}
\caption{Service level by model: Croston buys M5's best service (93.9\%);
Store service sits far lower for all models (43--74\%), reflecting denser,
higher-volume exposure. Same source as \cref{fig:totalcost}.}
\label{fig:service}
\end{figure}
\begin{figure}[H]
\centering
\includegraphics[width=0.95\textwidth]{03_cost_breakdown_stacked.png}
\caption{Holding vs.\ stockout composition. Naive's Store bar is almost entirely
stockout; DES/TES bars are holding-heavy everywhere --- two different ways to
waste money. Same source as \cref{fig:totalcost}.}
\label{fig:costsplit}
\end{figure}
\begin{figure}[H]
\centering
\includegraphics[width=0.95\textwidth]{04_inventory_levels_by_model.png}
\caption{Average inventory and reorder frequency: DES/TES carry roughly double
the stock of the leaders. Same source as \cref{fig:totalcost}.}
\label{fig:invlevels}
\end{figure}

\subsection{Forecast accuracy vs.\ inventory decision quality}
\label{sec:paradox}
\Cref{fig:scatter} is the report's center of gravity. Each point is a model:
horizontal position is forecast error, vertical is inventory cost. On M5 the
cloud slopes sensibly --- better forecasts, cheaper shelves, LSTM at the good
corner. On Store the relationship breaks: LSTM sits leftmost (best error) but
mid-pack vertically, while Moving Average, SES, and Seasonal Naive sit to its
right and \emph{below} it. The mechanism is cost asymmetry: the policy charges
under-prediction (lost sales at $P=5$) five times harder than over-prediction
(holding at $H=1$), and safety stock scales with each model's own error
volatility. A slightly biased-toward-averaging forecaster that rarely
under-shoots can out-stock a sharper but occasionally short one. Accuracy is a
symmetric ruler; the shelf is not symmetric.

\begin{figure}[H]
\centering
\includegraphics[width=0.9\textwidth]{05_forecast_vs_inventory_scatter.png}
\caption{Forecast error vs.\ inventory cost per model. M5: aligned (better
forecasts, cheaper shelves). Store: broken (LSTM leftmost but mid-cost).
Same source as \cref{fig:totalcost}.}
\label{fig:scatter}
\end{figure}

\subsection{Croston-family findings}
\label{sec:croston-find}
The specialists justify their rung. On M5, TSB is the family's best forecaster
(MAE 1.010, MASE 1.369 --- fourth overall), SBA second (1.066), Croston third
(1.082): adaptivity beats bias-correction beats the original on accuracy. On
cost, the order reverses its head: SBA (160.66) $<$ TSB (165.13) $<$ Croston
(174.81). SBA's debiasing trims exactly the over-forecast that holding cost
punishes, while TSB's responsiveness costs a little extra stock. Two lessons:
even inside one family, accuracy rank $\neq$ cost rank; and the family as a
whole validates the design choice --- on intermittent demand, models that
refuse to fit zeros belong in the comparison.

\begin{figure}[H]
\centering
\includegraphics[width=0.9\textwidth]{01_croston_intuition.png}
\caption{Croston intuition: split sparse demand into \emph{how much when it
happens} and \emph{how long between happenings}, smooth each separately.
Teaching figure from notebook 06c.}
\label{fig:croston-intuition}
\end{figure}
\begin{figure}[H]
\centering
\includegraphics[width=0.95\textwidth]{01_comparison_MAE_with_croston.png}
\caption{M5 MAE with the Croston family: TSB leads the family at 1.010.
Source: notebook 06c.}
\label{fig:croston-mae}
\end{figure}

\subsection{Demand-archetype findings}
\label{sec:arch-find}
\Cref{fig:arch-mae,fig:arch-wape} break M5 accuracy by demand shape (pooled
WAPE/sMAPE; per-series ratios would let one near-zero series dominate).
LSTM wins Highly Intermittent (MAE 0.416) and Intermittent (0.989) outright.
In Variable demand the top three --- SES (1.799), Moving Average (1.804), LSTM
(1.806) --- sit within 0.007 of each other: a genuine three-way tie, honestly
reported as such. SES takes Moderate (2.634). TES is worst in \emph{every}
archetype, including the smooth ones its trend-seasonality machinery should
favor. The Smooth column (single series; LSTM nominally best at forecast MAE
2.358) is anecdotal and carries no conclusion --- it is shown for completeness,
not evidence. Inventory-by-archetype (\cref{fig:arch-cost}) splits the crown:
LSTM is cheapest in Highly Intermittent (64.52) and Intermittent (147.94),
Moving Average takes Variable (285.93), and SBA takes Moderate (383.59) ---
the intermittent-demand specialist winning exactly where intermittency bites
least smoothly. Sparse structure rewards the learner; but no single method
owns every demand shape, which is itself a finding.

\begin{figure}[H]
\centering
\includegraphics[width=0.95\textwidth]{mae_by_archetype.png}
\caption{M5 MAE by archetype: LSTM leads sparse shapes; Variable is a
three-way tie; TES trails everywhere. Pooled-metric figure (notebook
archetype layer).}
\label{fig:arch-mae}
\end{figure}
\begin{figure}[H]
\centering
\includegraphics[width=0.95\textwidth]{wape_by_archetype.png}
\caption{M5 pooled WAPE by archetype: the ratio-of-sums construction keeps
near-zero series from exploding the mean. Same source as
\cref{fig:arch-mae}.}
\label{fig:arch-wape}
\end{figure}
\begin{figure}[H]
\centering
\includegraphics[width=0.9\textwidth]{03_arima_by_archetype_m5.png}
\caption{ARIMA error by M5 archetype: degradation concentrates where demand is
sparsest --- the visual companion to the 22.65\% fallback rate. Source:
notebook 07.}
\label{fig:arima-arch}
\end{figure}
\begin{figure}[H]
\centering
\includegraphics[width=0.9\textwidth]{06_cost_by_archetype_m5.png}
\caption{M5 inventory cost by archetype: LSTM cheapest on sparse shapes;
Moving Average takes Variable, SBA takes Moderate. Same source as \cref{fig:totalcost}.}
\label{fig:arch-cost}
\end{figure}

\subsection{LSTM findings}
\label{sec:lstm-find}
The neural rung half-validates its thesis. As a forecaster the global LSTM is
untouchable: best MAE and MASE on both datasets, best archetype MAE on sparse
shapes, sub-1.0 MASE on Store (the only model to beat the training benchmark).
As an inventory driver it splits: best on M5 (152.83, robust at 25/27) but
fourth on Store (2,247.45), where its sharper forecasts carry an
under-prediction profile the $P/H=5$ cost ratio punishes. The hardening note
matters here: the LSTM was revalidated on chronological per-series splits
after overlapping-window leakage was found in the first validation design, so
these numbers reflect genuine generalization, not validation peeking.
Bottom line for practitioners: buy the LSTM for forecast quality with
confidence; for shelf decisions, check your cost ratio first.

\subsection{ARIMA findings}
\label{sec:arima-find}
ARIMA is the fragile linear middle. On dense Store data it fits cleanly
(4,000/4,000, zero fallbacks) yet finishes sixth on both boards (MASE 1.452;
cost 2,320.87) --- linear memory adds little where weekly structure dominates.
On sparse M5 data, 906 of 4,000 fits (22.65\%) trip the $>85\%$-zero pre-check
and defer to Naive, concentrated in Highly Intermittent series
(\cref{fig:arima-arch}); crucially, there are zero optimizer failures ---
every non-skipped fit converges. The pre-check is a designed guardrail, not a
crash log, and its asymmetry (22.65\% vs.\ 0\%) is itself a measurement of how
demand structure dictates statistical feasibility.

\subsection{SARIMA full-population findings}
\label{sec:sarima-find}
The seasonal challenger delivers. All 4,000 Store fits succeed with zero
fallbacks in 2,382\,s ($\approx$0.6\,s/fit): seasonal differencing at $m=7$ is
feasible at full population, not just on the exploratory 100-series subset
(MAE 8.31) that preceded it. Final full500 accuracy (MAE 8.45, MASE 1.055,
RMSSE 1.040) shifts slightly from the subset --- a population effect, logged
openly. Decision quality (2,282.31, fifth) lags accuracy (second): strong
weekly tracking still under-protects against stockouts at $P/H=5$. SARIMA was
never run on M5, deliberately: seasonal differencing of mostly-zero series is
methodologically unsound, and the Croston rung already covers that territory.

\subsection{Model complexity findings}
\label{sec:complexity}
H3 in monotone form is rejected. The most complex classical models are the
worst performers (TES last everywhere; DES an overstocking liability), the
mid-complexity linear model sits mid-pack, and only the neural model converts
complexity into accuracy --- and even then into inventory value on just one of
two datasets. Meanwhile the simplest rung (Moving Average, SES) wins or
near-wins three of the four boards. Complexity buys \emph{capability}
(LSTM's cross-series learning is real), not \emph{value}: value is decided by
the interaction of capability with demand structure and cost asymmetry.

\subsection{Robustness across inventory policies}
\label{sec:robust}
\Cref{tab:robust} with \cref{fig:heat-m5,fig:heat-store} is the trust map.
M5: LSTM first in 25 of 27 policies --- robust. Its only defeats, both to SES,
sit in one corner: short lead time ($L=3$) with high stockout penalty
($P=10$), where SES's conservative level-tracking wastes less on emergency
cover. Store: fragile --- Moving Average 9, SES 7, Naive 6, DES 5. No model
reaches even half the grid; the ``winner'' is a property of the policy, not
the model. This asymmetry is the report's second major device after the
forecast--decision gap: bank M5-LSTM; rent Store-MA.

\begin{table}[H]
\centering
\caption{Robustness verdict: policy wins across the 27-policy grid. M5 robust; Store fragile.}
\label{tab:robust}
\begin{tabular}{lll}
\toprule
Dataset & Wins (of 27) & Verdict \\
\midrule
M5 & LSTM 25, SES 2 (both at $L{=}3$, $P{=}10$) & \textbf{Robust} \\
Store & MA 9, SES 7, Naive 6, DES 5 & \textbf{Fragile / policy-sensitive} \\
\bottomrule
\end{tabular}
\end{table}

\begin{figure}[H]
\centering
\includegraphics[width=0.9\textwidth]{rank_heatmap_m5.png}
\caption{M5 rank heatmap across 27 policies: an unbroken LSTM band with two
SES cells in the short-lead/high-penalty corner. Source: sensitivity layer.}
\label{fig:heat-m5}
\end{figure}
\begin{figure}[H]
\centering
\includegraphics[width=0.9\textwidth]{rank_heatmap_store_item_demand.png}
\caption{Store rank heatmap: leadership fragments across four models as the
policy moves. Same source as \cref{fig:heat-m5}.}
\label{fig:heat-store}
\end{figure}
\begin{figure}[H]
\centering
\includegraphics[width=0.9\textwidth]{F-NEW-03-sensitivity-winner-share.png}
\caption{Winner shares across the 27-policy grid: M5 concentration vs.\ Store
fragmentation. Deterministic rebuild (F-NEW-03).}
\label{fig:winnershare}
\end{figure}

\subsection{Sensitivity analysis}
\label{sec:sens}
The grid's interior tells managers \emph{when} leadership changes. Longer lead
times amplify safety-stock differences (volatility-driven models suffer);
higher $P/H$ punishes under-predictors (Naive collapses everywhere except
short-lead Store corners where its leanness pays); higher service targets
reward conservative levels (SES gains). {\sloppy Per-archetype M5 costing (55 rows,
\texttt{sensitivity\_by\_archetype\_m5.csv}) confirms the structure story:
sparse archetypes magnify every policy effect.\par} Full matrices are
\cref{app:sens}.

\subsection{Convergence and feasibility}
\label{sec:conv}
\Cref{tab:conv} accounts for every fit. M5 ARIMA: 3,094 fits, 906 pre-check
skips (22.65\%), 0 optimizer failures. Store ARIMA: 4,000/4,000. SARIMA
subset (exploratory): 800/800. SARIMA full500 (primary): 4,000/4,000 in
2,382\,s. The feasibility gradient --- dense fits always, sparse fits
sometimes --- is a result about the world, not about the software.

\begin{table}[H]
\centering
\caption{Convergence accounting. Pre-check skips are by-design guards, not failures.}
\label{tab:conv}
\begin{tabular}{llrrrr}
\toprule
Scope & Model & Attempted & Failed & Naive fallback & Fallback \% \\
\midrule
M5 full500 & ARIMA & 4,000 & 0 & 906 & 22.65 \\
Store full500 & ARIMA & 4,000 & 0 & 0 & 0.00 \\
Store subset (explor.) & SARIMA & 800 & 0 & 0 & 0.00 \\
Store full500 (primary) & SARIMA & 4,000 & 0 & 0 & 0.00 \\
\bottomrule
\end{tabular}
\end{table}

\subsection{Negative and surprising findings}
\label{sec:negative}
Stated plainly, because nulls are findings: (1) TES is worst on all four
boards and every archetype. (2) DES buys double inventory for worse service
than SES. (3) Naive --- the dumbest method --- wins 6 Store policies. (4)
Seasonal Naive is Store's inventory bronze medalist. (5) The LSTM loses Store
inventory to a moving average. (6) An intermediate WAPE computation exploded to
$\sim 2\times 10^9$ on near-zero-test-demand series before pooling fixed it
--- a live demonstration of why metric construction needs the same skepticism
as modeling. (7) MA vs.\ SES on M5 is statistically indistinguishable in
practice despite $n=112{,}000$.

\subsection{Consolidated findings and hypothesis verdicts}
\label{sec:verdicts}
\textbf{H1 (accuracy differs): confirmed} --- 91 paired tests, nearly all
significant with Holm control. \textbf{H2 (accuracy moves inventory):
confirmed with the central qualification} --- it moves inventory, but not
monotonically and not always in the same direction (Store paradox).
\textbf{H3 (monotone sophistication): rejected} --- TES/DES/ARIMA lose to
simpler methods; only LSTM complexity pays, and only partially. \textbf{H4
(cross-dataset consistency): rejected in its strong form} --- forecasting
leadership transfers (LSTM both), decision leadership does not (LSTM vs.\ MA),
and robustness transfers not at all (robust vs.\ fragile). The qualified
verdict is the contribution: \emph{what} transfers (forecast skill) and
\emph{what does not} (decision superiority, policy-stability) are now mapped.
